% ============================================================================
%  Appendix -- the tanh front, its width, and its gradient integral
%
%  This is the derivation the draft has referred to since v3 without ever
%  containing it.  In v6 Appendix E is the parameter table, so this becomes
%  whichever letter follows it.
%
%  It earns its place three times over: it supplies w = sqrt(2 kappa), which is
%  the gate of the validity section; it supplies I = 4/(3w), which is what closes
%  the speed formula; and the travelling-frame version at the end explains, in
%  two lines, why the closed form overpredicts at low damping and underpredicts
%  at high drive -- the sign flip that would otherwise be an unexplained
%  empirical caveat.
%
%  The two width corrections are named and their direction explained, but NOT
%  quantified: the fitted form w = sqrt(2(kappa - v^2))(1 + 1.53 Pi^2) is
%  deliberately left out of this write-up.  It is overdamped-only (it inverts at
%  Omega = 1), its coefficient is fitted, and it still leaves 4.8% worst case --
%  more to defend than it is worth here.  The mechanisms cost nothing and are
%  what answer "why does the error change sign".
%
%  provenance: results/closed_form_validation.csv for every measured number.
% ============================================================================

\section{The tanh front and its gradient integral}
\label{app:tanh}

\subsection*{The continuum limit}

Section \ref{sec:boyle} wrote the gas force exactly as
%
\begin{equation}
  F_n = \eta\theta \,\bigl(\tilde u_{n+1} - 2\tilde u_n + \tilde u_{n-1}\bigr)
        \frac{1}{\tilde V_{n-1}\tilde V_n},
\end{equation}
%
a spring-chain force of stiffness $\kappa = \eta\theta$ times a Boyle correction.
For a front spread over many lattice sites the strains $\Delta_n$ are of order
$1/w$, so the correction tends to unity, and the discrete Laplacian tends to a
second derivative with respect to the continuous coordinate $\xi$. Writing
$\tilde u_n(\tilde t) = U(\xi)$ the equation of motion becomes
%
\begin{equation}
  \kappa\,U'' = U^3 - U + \Pi,
  \label{eq:static-kink}
\end{equation}
%
in the static, undriven limit --- the two corrections to which are taken up at
the end of this appendix.

\subsection*{The first integral}

Equation \eqref{eq:static-kink} at $\Pi = 0$ integrates once. Multiplying through
by $U'$ turns both sides into exact derivatives,
%
\begin{equation}
  \kappa\,U'U'' = \bigl(U^3 - U\bigr)U'
  \quad\Longrightarrow\quad
  \frac{\mathrm{d}}{\mathrm{d}\xi}\Bigl[\tfrac12 \kappa (U')^2\Bigr]
  = \frac{\mathrm{d}}{\mathrm{d}\xi}\Bigl[\tfrac14 U^4 - \tfrac12 U^2\Bigr],
\end{equation}
%
so that
%
\begin{equation}
  \tfrac12 \kappa (U')^2 = \tfrac14 U^4 - \tfrac12 U^2 + C.
  \label{eq:first-integral}
\end{equation}
%
The boundary conditions fix $C$, and it is worth being explicit about why they
take the form they do. The front connects the two stable states, so
$U \to \pm 1$ as $\xi \to \pm\infty$. And the material far ahead of and far
behind the front is sitting still in a well rather than in the middle of
switching, so $U' \to 0$ at both ends. It is that vanishing gradient which kills
the integration constant: substituting $U = 1$, $U' = 0$ into
\eqref{eq:first-integral} gives $C = \tfrac14$, and hence
%
\begin{equation}
  \tfrac12 \kappa (U')^2 = \tfrac14\bigl(U^2 - 1\bigr)^2,
  \qquad
  U' = \frac{1 - U^2}{\sqrt{2\kappa}} ,
\end{equation}
%
taking the branch with $U' > 0$ for $|U| < 1$, which is the front that switches
upward. Had $C$ been anything else the right-hand side would not have been a
perfect square and the solution would have been periodic rather than a single
kink.

Separating variables,
%
\begin{equation}
  \int \frac{\mathrm{d}U}{1 - U^2} = \int \frac{\mathrm{d}\xi}{\sqrt{2\kappa}}
  \quad\Longrightarrow\quad
  \operatorname{artanh} U = \frac{\xi - \xi_0}{\sqrt{2\kappa}},
\end{equation}
%
and inverting gives the front,
%
\begin{equation}
  \boxed{\;U(\xi) = \tanh\!\left(\frac{\xi - \xi_0}{\sqrt{2\kappa}}\right),
  \qquad w = \sqrt{2\kappa}\;}
  \label{eq:tanh}
\end{equation}
%
with $\xi_0$ free, as it must be: the lattice equation is translation invariant,
and it is precisely that freedom which the Peierls--Nabarro barrier destroys once
the front is narrow enough to feel individual sites.

\subsection*{The gradient integral}

The energy-balance speed of Section \ref{sec:speed} needs
$I = \int_{-\infty}^{\infty} (U')^2 \,\mathrm{d}\xi$. From \eqref{eq:tanh},
$U' = w^{-1}\operatorname{sech}^2(\xi/w)$, and with
$\int \operatorname{sech}^4 x \,\mathrm{d}x = 4/3$ over the whole line,
%
\begin{equation}
  I = \frac{1}{w^{2}}\int_{-\infty}^{\infty}
      \operatorname{sech}^{4}\!\left(\frac{\xi}{w}\right) \mathrm{d}\xi
    = \frac{1}{w^{2}} \cdot w \cdot \frac{4}{3}
    = \frac{4}{3w}.
  \label{eq:I}
\end{equation}
%
Together with $\Delta\tilde\psi = 2\Pi$ from Appendix \ref{app:dpsi} this closes
the speed formula,
%
\begin{equation}
  v_{\rm pred} = \frac{\Delta\tilde\psi}{\Omega I}
         = \frac{2\Pi}{\Omega}\cdot\frac{3w}{4}
         = \frac{3\,\Pi\sqrt{2\kappa}}{2\,\Omega}.
  \label{eq:closed}
\end{equation}

\subsection*{Two corrections, and the sign flip}

Equation \eqref{eq:tanh} is the \emph{static, undriven} kink, and each of the two
assumptions behind it fails in a direction we can now name.

\paragraph{Motion narrows the front.} Restoring the time dependence with
$\tilde u_n = U(n - v\tilde t)$, so that
$\partial_{\tilde t}^2 \to v^2 U''$ and $\partial_{\tilde t} \to -vU'$, the
equation of motion becomes
%
\begin{equation}
  \bigl(\kappa - v^{2}\bigr) U'' + \Omega v\, U' = U^3 - U + \Pi .
  \label{eq:travelling}
\end{equation}
%
The stiffness that sets the width is therefore not $\kappa$ but $\kappa - v^2$: a
moving front is narrower than a static one, and the effect is largest exactly
where $v$ is largest, which is at low damping. (Equation \eqref{eq:travelling}
also shows there can be no monotone kink at $v^2 > \kappa$.)

\paragraph{Drive broadens it.} At $\Pi \neq 0$ the well tilts, the minima move to
$\tilde u_\pm = \pm 1 - \Pi/2 + O(\Pi^2)$, and the front is no longer the
zero-drive $\tanh$. The correction must be \emph{even} in $\Pi$, since sending
$\Pi \to -\Pi$ maps the potential to its mirror image and the kink to its mirror
image, which has the same width; the leading term is therefore $O(\Pi^2)$. It is
not negligible at the top of the range sampled here: at $\Pi = 0.3$ the measured
front is more than $10\%$ broader than $\sqrt{2\kappa}$.

\paragraph{Why the closed form flips sign.} The two corrections push the width
opposite ways, and \eqref{eq:closed} is linear in $w$, so they push the predicted
speed opposite ways too. At low damping the front runs fast, $\kappa - v^2$ bites,
the true front is narrower than $\sqrt{2\kappa}$ and \eqref{eq:closed}
\emph{over}predicts --- by up to $12.7\%$ at $\Omega = 1$. At high damping the
front is slow, that term is negligible, the $\Pi$ tilt dominates, the true front
is broader and \eqref{eq:closed} \emph{under}predicts --- by up to $10.2\%$ at
$\Omega \ge 3$. Between them the two cancel, which is why the closed form is at
its most accurate near $\Omega \approx 1$--$3$ and why its useful window is
stated as a range in $\Omega$ rather than a bound.
